\subsubsection{Data Management}
\begin{enumerate}
    \item \textbf{Data Reading and processing}
          
          \begin{itemize}
            \item \textbf{Read Gene IDs from TSV File (R)}
                  
                  The function reads gene identifiers from a tab-separated value (TSV) file,
                  extracting specific gene IDs from a designated column while skipping
                  header rows. This function is typically used as the first step in processing
                  gene expression data to obtain the complete list of gene identifiers. Call
                  the \texttt{read\_gene\_ids\_from\_tsv\_file} function with the following arguments:
                  
                  \begin{itemize}
                    \item \texttt{filename}: Character string of the filename
                          
                    \item \texttt{ngenes}: Number of gene IDs to read
                          
                    \item \texttt{gene\_len}: Maximum length of each gene ID
                          
                    \item \textbf{(Optional)} \texttt{n\_header\_rows}: Number of header
                          rows to skip
                          
                    \item \textbf{(Optional)} \texttt{gene\_col}: Column index (1-based) of
                          gene IDs in the file
                  \end{itemize}
                  
                  \textbf{Return list:}
                  \begin{itemize}
                    \item \textbf{gene\_ids}: Character vector of gene IDs read from the file
                          
                    \item \textbf{ierr}: Integer error code (0 if successful)
                  \end{itemize}
                  
                  \begin{lstlisting}[language=R, caption={R Read Gene IDs from TSV File}]
read_gene_ids_from_tsv_file(filename, ngenes, gene_len, 
    n_header_rows = 1, gene_col = 1)
                  \end{lstlisting}
                  
            \item \textbf{Read Expression Vectors (R)}
                  
                  The function reads gene expression data from multiple tabular files (CSV/TSV
                  format), extracting expression values for specified genes across multiple
                  samples. This function consolidates expression data from multiple files into
                  a single matrix. Call the \texttt{read\_expression\_vectors} function with
                  the following arguments:
                  
                  \begin{itemize}
                    \item \texttt{file\_list}: Character vector of filenames
                          
                    \item \texttt{gene\_ids}: Character vector of gene IDs to match
                          
                    \item \texttt{n\_header\_rows}: Number of header rows to skip in each file
                          
                    \item \texttt{gene\_col}: Column index (1-based) of gene IDs in each file
                          
                    \item \texttt{value\_cols}: Vector of column indices (1-based) of expression
                          values in each file
                          
                    \item \texttt{n\_samples}: Total number of samples (rows) to read across
                          all files
                          
                    \item \textbf{(Optional)} \texttt{delimiter}: Delimiter used in the
                          files (default: tab)
                  \end{itemize}
                  
                  \textbf{Return list:}
                  \begin{itemize}
                    \item \textbf{expression\_vectors}: Numeric matrix of dimensions (n\_samples
                          × n\_genes)
                          
                    \item \textbf{ierr}: Integer error code (0 if successful)
                  \end{itemize}
                  
                  \begin{lstlisting}[language=R, caption={R Read Expression Vectors}]
read_expression_vectors(file_list, gene_ids, n_header_rows, gene_col, 
    value_cols, n_samples, delimiter = "\t")
                  \end{lstlisting}
                  
            \item \textbf{Read OrthoFinder File (R)}
                  
                  The function processes an OrthoFinder output file to map genes to their respective
                  orthologous families. This function establishes the gene-family
                  relationships essential for downstream orthology-based analyses. Call the \texttt{read\_orthofinder\_file}
                  function with the following arguments:
                  
                  \begin{itemize}
                    \item \texttt{filename}: Character string of the filename
                          
                    \item \texttt{gene\_ids}: Character vector of gene IDs to match
                          
                    \item \texttt{n\_families}: Number of unique gene families expected
                          
                    \item \texttt{family\_len}: Maximum length of each family ID
                  \end{itemize}
                  
                  \textbf{Return list:}
                  \begin{itemize}
                    \item \textbf{family\_ids}: Character vector of family IDs read from the
                          file
                          
                    \item \textbf{gene\_to\_fam}: Integer vector mapping each gene to its family
                          index (0 = unassigned)
                          
                    \item \textbf{ierr}: Integer error code (0 if successful)
                  \end{itemize}
                  
                  \textbf{Note:} If genes are not found in the gene IDs list a message will
                  be printed but NO error code is thrown.
                  
                  \begin{lstlisting}[language=R, caption={R Read OrthoFinder File}]
read_orthofinder_file(filename, gene_ids, n_families, family_len)
                  \end{lstlisting}
                  
                  \textbf{Complete R Data Loading Workflow Example:}
                  \begin{lstlisting}[language=R, caption={R Complete Data Loading Workflow}]
# Step 1: Load gene identifiers
gene_result <- read_gene_ids_from_tsv_file(
    filename = "data/expression_data.tsv",
    ngenes = 50000,
    gene_len = 32,
    n_header_rows = 1,
    gene_col = 1
)
gene_ids <- gene_result$gene_ids

# Step 2: Load expression data from multiple samples
samples <- c("sample1.tsv", "sample2.tsv", 
    "sample3.tsv", "sample4.tsv")

expression_result <- read_expression_vectors(
    file_list = samples,
    gene_ids = gene_ids,
    n_header_rows = 1,
    gene_col = 1,
    value_cols = 2,
    n_samples = 4
)
expression_vectors <- expression_result$expression_vectors

# Step 3: Load orthology information
family_result <- read_orthofinder_file(
    filename = "data/Orthogroups.tsv",
    gene_ids = gene_ids,
    n_families = 15000,
    family_len = 32
)
family_ids <- family_result$family_ids
gene_to_fam <- family_result$gene_to_fam
                  \end{lstlisting}
          \end{itemize}
    \item \textbf{Data Filtering and Validation}
          \begin{itemize}
            \item \textbf{Filter Unassigned Genes (R)}
                  
                  The function filters out genes that are not assigned to any family (where
                  gene\_to\_fam = 0). This is essential for cleaning data before orthology-based
                  analyses, ensuring only genes with valid family assignments are processed.
                  Call the \texttt{filter\_unassigned\_genes} function with the following arguments:
                  
                  \begin{itemize}
                    \item \texttt{gene\_ids}: Character vector of gene IDs
                          
                    \item \texttt{expression\_vectors}: Numeric matrix of expression values (n\_samples
                          × n\_genes)
                          
                    \item \texttt{gene\_to\_fam}: Integer vector mapping each gene to its family
                          index (0 = unassigned)
                  \end{itemize}
                  
                  \textbf{Return list:}
                  \begin{itemize}
                    \item \textbf{gene\_ids}: Filtered character vector of gene IDs
                          
                    \item \textbf{expression\_vectors}: Filtered numeric matrix of expression
                          values
                          
                    \item \textbf{gene\_to\_fam}: Filtered integer vector mapping each gene to
                          its family index
                          
                    \item \textbf{n\_genes\_kept}: Number of genes kept after filtering
                  \end{itemize}
                  
                  \begin{lstlisting}[language=R, caption={R Filter Unassigned Genes}]
filter_unassigned_genes(gene_ids, expression_vectors, gene_to_fam)
                  \end{lstlisting}
                  
            \item \textbf{Validate Gene to Family Mapping (R)}
                  
                  The function validates the gene-to-family mapping array, ensuring all
                  family indices are within valid range and properly assigned. This prevents
                  downstream errors from invalid family indices. Call the \texttt{validate\_gene\_to\_family\_mapping}
                  function with the following arguments:
                  
                  \begin{itemize}
                    \item \texttt{gene\_to\_fam}: Integer vector mapping each gene to its family
                          index (0 = unassigned)
                          
                    \item \texttt{n\_genes}: Number of genes
                          
                    \item \texttt{n\_families}: Number of gene families
                  \end{itemize}
                  
                  \begin{lstlisting}[language=R, caption={R Validate Gene to Family Mapping}]
validate_gene_to_family_mapping(gene_to_fam, n_genes, n_families)
                  \end{lstlisting}
                  
            \item \textbf{Validate Expression Data (R)}
                  
                  The function validates expression data, checking for NaN/Inf values and
                  optionally verifying non-negative values. This ensures data quality before
                  computational analyses. Call the \texttt{validate\_expression\_data} function
                  with the following arguments:
                  
                  \begin{itemize}
                    \item \texttt{expression\_vectors}: Numeric matrix of expression values (n\_samples
                          × n\_genes)
                          
                    \item \texttt{n\_genes}: Number of genes
                          
                    \item \texttt{n\_samples}: Number of samples
                          
                    \item \textbf{(Optional)} \texttt{check\_non\_negative}: Logical flag to
                          check for non-negative values (default: TRUE)
                  \end{itemize}
                  
                  \begin{lstlisting}[language=R, caption={R Validate Expression Data}]
validate_expression_data(expression_vectors, n_genes, n_samples, 
    check_non_negative = TRUE)
                  \end{lstlisting}
                  
            \item \textbf{Validate Family Centroids (R)}
                  
                  The function validates family centroids data, checking for NaN/Inf values
                  and proper data structure. Centroids represent the central tendency of each
                  gene family's expression profile. Call the \texttt{validate\_family\_centroids}
                  function with the following arguments:
                  
                  \begin{itemize}
                    \item \texttt{family\_centroids}: Numeric matrix of family centroids (n\_samples
                          × n\_families)
                          
                    \item \texttt{n\_families}: Number of gene families
                          
                    \item \texttt{n\_samples}: Number of samples
                  \end{itemize}
                  
                  \begin{lstlisting}[language=R, caption={R Validate Family Centroids}]
validate_family_centroids(family_centroids, n_families, n_samples)
                  \end{lstlisting}
                  
            \item \textbf{Validate Shift Vectors (R)}
                  
                  The function validates shift vectors, ensuring data types are correct and
                  structure matches expected patterns. Shift vectors represent the deviation
                  of each gene from its family centroid. Call the \texttt{validate\_shift\_vectors}
                  function with the following arguments:
                  
                  \begin{itemize}
                    \item \texttt{shift\_vectors}: Numeric matrix of shift vectors (2*n\_samples
                          × n\_genes)
                          
                    \item \texttt{expression\_vectors}: Numeric matrix of expression values (n\_samples
                          × n\_genes)
                          
                    \item \texttt{family\_centroids}: Numeric matrix of family centroids (n\_samples
                          × n\_families)
                          
                    \item \texttt{gene\_to\_fam}: Integer vector mapping each gene to its family
                          index
                          
                    \item \texttt{n\_samples}: Number of samples
                  \end{itemize}
                  
                  \begin{lstlisting}[language=R, caption={R Validate Shift Vectors}]
validate_shift_vectors(shift_vectors, expression_vectors, 
    family_centroids, gene_to_fam, n_samples)
                  \end{lstlisting}
                  
            \item \textbf{Validate String Array Uniqueness (R)}
                  
                  The function validates that all strings in an array are unique, using a
                  hashset internally. Note: This may increase memory usage temporarily for
                  large datasets. Call the \texttt{validate\_string\_array\_uniqueness}
                  function with the following arguments:
                  
                  \begin{itemize}
                    \item \texttt{string\_arr}: Character vector of gene IDs
                          
                    \item \texttt{n\_strings}: Number of genes
                  \end{itemize}
                  
                  \begin{lstlisting}[language=R, caption={R Validate String Array Uniqueness}]
validate_string_array_uniqueness(string_arr, n_strings)
                  \end{lstlisting}
                  
            \item \textbf{Validate All Data (R)}
                  
                  The function performs comprehensive validation of all data components in one
                  call. This function executes all individual validations sequentially for
                  complete data quality assurance. Call the \texttt{validate\_all\_data}
                  function with the following arguments:
                  
                  \begin{itemize}
                    \item \texttt{n\_genes}: Number of genes
                          
                    \item \texttt{n\_families}: Number of gene families
                          
                    \item \texttt{n\_samples}: Number of samples
                          
                    \item \texttt{gene\_ids}: Character vector of gene IDs
                          
                    \item \texttt{gene\_family\_ids}: Character vector of gene family IDs
                          
                    \item \texttt{gene\_to\_fam}: Integer vector mapping each gene to its family
                          index (0 = unassigned)
                          
                    \item \texttt{expression\_vectors}: Numeric matrix of expression values (n\_samples
                          × n\_genes)
                          
                    \item \texttt{family\_centroids}: Numeric matrix of family centroids (n\_samples
                          × n\_families)
                          
                    \item \texttt{shift\_vectors}: Numeric matrix of shift vectors (2*n\_samples
                          × n\_genes)
                  \end{itemize}
                  
                  \begin{lstlisting}[language=R, caption={R Validate All Data}]
validate_all_data(n_genes, n_families, n_samples, gene_ids, 
    gene_family_ids, gene_to_fam, expression_vectors, 
    family_centroids, shift_vectors)
                  \end{lstlisting}
                  
                  \textbf{R Data Filtering and Validation Workflow Example:} 
                  \begin{lstlisting}[language=R, caption={R Data Filtering and Validation Workflow}]
# After loading data, filter unassigned genes
filter_result <- filter_unassigned_genes(gene_ids, expression_vectors, 
    gene_to_fam)

gene_ids <- filter_result$gene_ids
expression_vectors <- filter_result$expression_vectors
gene_to_fam <- filter_result$gene_to_fam
n_genes_kept <- filter_result$n_genes_kept

# Run comprehensive validations
validate_string_array_uniqueness(gene_ids, n_genes_kept)
validate_string_array_uniqueness(family_ids, n_families)
validate_expression_data(expression_vectors, n_genes_kept, 
    n_samples, check_non_negative = TRUE)

validate_gene_to_family_mapping(gene_to_fam, n_genes_kept, n_families)

# For complete analysis with centroids and shift vectors
validate_all_data(
    n_genes_kept, n_families, n_samples, 
    gene_ids, family_ids, gene_to_fam,
    expression_vectors, family_centroids, shift_vectors
)
            \end{lstlisting}
          \end{itemize}
\end{enumerate}